%!TEX root = ../main.tex

\addtocontents{toc}{\protect\vskip1em}
\chapter*{Conclusioni}\label{ch:conclusioni}
\markboth{Conclusioni}{}
\addcontentsline{toc}{chapter}{Conclusioni}

I materiali raccolti in questo documento permettono di impostare la relazione finale di tirocinio in modo ordinato e coerente con il progetto effettivamente analizzato. La raccolta delle informazioni evidenzia che \builder{} non e soltanto un editor grafico, ma un'applicazione web articolata che integra modellazione ER, trasformazione verso lo schema logico, esportazione, salvataggio, localizzazione, reverse engineering SQL e test.

La struttura proposta per la relazione consente di mantenere il giusto equilibrio tra descrizione tecnica e dimensione formativa del tirocinio. In particolare, la relazione finale dovra valorizzare il ruolo della Prof.ssa Alessandra Lumini come tutor, sottolineando la sua supervisione e il contributo dei suoi suggerimenti nel definire priorita, chiarezza, qualita, usabilita e correttezza concettuale del lavoro.

Prima della stesura definitiva restano da completare alcuni dati amministrativi, come periodo del tirocinio, numero di ore, CFU, corso di laurea, struttura ospitante, dati dello studente ed eventuali screenshot da inserire. Una volta completate queste informazioni, il prompt incluso nel documento potra essere usato in Prism per generare una prima bozza della relazione finale.
